\chapter{Small Practical Programs}\label{ch:projects}

\begin{goals}
\begin{itemize}
  \item Bringing everything you have learned together in a few complete programs
  \item Breaking an everyday problem into small steps
  \item Reading a program with several functions and following its flow
  \item Ideas for the road ahead
\end{itemize}
\end{goals}

\section{From tools to programs}

In the earlier chapters you met the basic tools of programming one at a
time: printing, variables, types, arithmetic, conditions, loops, functions,
arrays and strings. Every program ever written, however large, is built from
these same tools. In this chapter we write a few small but complete
programs; programs that resemble real work, and in each of which several
tools work together.

For each program we first state the problem, then break it into small steps,
and then write the program. We suggest that before reading our program you
try to write it yourself, and then compare the two versions. There is
nothing wrong with your way being different from ours.

\section{An electricity bill}

\textbf{The problem.} The electricity company charges for monthly
consumption in tiers: the first 100 kilowatt-hours cost 5 cents each, the
next 200 units cost 10 cents each, and anything beyond that costs 20 cents
a unit. Work out the bills of several customers.

\textbf{The steps.} We want a function that takes the consumption and
returns the amount. We use the consumption up tier by tier: first up to 100
units at the first rate, then up to 200 units at the second rate, and the
rest at the third. Then, for an array of consumptions, we print each
customer's bill.

\program{A tiered electricity bill}
\begin{salamcode}
func bill amount(usage: int): int:
    mut remaining := usage
    mut amount := 0

    // first tier: up to 100 units, 5 cents each
    tier one := remaining > 100 ? 100 : remaining
    amount += tier one * 5
    remaining -= tier one

    // second tier: up to the next 200 units, 10 cents each
    tier two := remaining > 200 ? 200 : remaining
    amount += tier two * 10
    remaining -= tier two

    // third tier: whatever is left, 20 cents each
    amount += remaining * 20
    ret amount
end

func main:
    customer usage := [80, 100, 250, 400]
    each usage in customer usage:
        bill := bill amount(usage)
        println "Usage", usage, "units:", bill, "cents"
    end
end
\end{salamcode}

\begin{salamout}
Usage 80 units: 400 cents
Usage 100 units: 500 cents
Usage 250 units: 2000 cents
Usage 400 units: 4500 cents
\end{salamout}

Follow the consumption of 250 units: 100 units in the first tier (500
cents), 150 units in the second tier (1500 cents), and nothing is left for
the third; 2000 cents in all. In each tier, the \code{? :} operator says
``the whole capacity of the tier, or whatever is left, whichever is
smaller''.

\begin{tip}[An idea to take further]
Turn the rates into \code{const} values and add a fourth tier. Then, in
\code{main}, work out the company's total income as well.
\end{tip}

\section{A chart of grades}

\textbf{The problem.} We have the grades of a class, and we want to see how
many people fall into each band (below 10, 10 to 14, 15 to 17, 18 to 20) and
show this as a chart made of stars.

\textbf{The steps.} We need four counters; for each grade we decide which
counter to increase. Then, for each band, we print one line with as many
stars as the counter says. Printing the stars is handed to a function.

\program{A star chart of grades}
\begin{salamcode}
func chart row(label: str, count: int):
    mut stars := ""
    repeat count:
        stars += "*"
    end
    println label + ":", stars, "(" + count + ")"
end

func main:
    grades := [12, 19, 8, 15, 20, 17, 11, 9, 16, 18, 14, 13]
    mut failed := 0
    mut average := 0
    mut good := 0
    mut excellent := 0

    each grade in grades:
        if grade < 10: failed++
        else grade < 15: average++
        else grade < 18: good++
        else: excellent++
        end
    end

    println "Chart of the class grades"
    chart row("below 10", failed)
    chart row("10 to 14", average)
    chart row("15 to 17", good)
    chart row("18 to 20", excellent)
end
\end{salamcode}

\begin{salamout}
Chart of the class grades
below 10: ** (2)
10 to 14: **** (4)
15 to 17: *** (3)
18 to 20: *** (3)
\end{salamout}

Notice the compact form of the chain of conditions: each branch has one
short instruction written on the same line, and a single \code{end} closes
the whole chain.

\begin{tip}[An idea to take further]
Instead of four separate counters, make an array of four zeros and increase
the right counter through its index. Then print the chart with one
\code{each} loop and a parallel array of labels.
\end{tip}

\section{A shop receipt}

\textbf{The problem.} We have the list of items in a customer's basket:
name, unit price and quantity. We must print a receipt that shows the amount
of each line, the total, the discount (10 percent if the purchase is over
500), and the final amount.

\textbf{The steps.} We have three parallel arrays. With one loop over the
names we work out and print the amount of each line and add it to the
total. After the loop we decide the discount with a condition.

\program{A shop receipt}
\begin{salamcode}
const threshold := 500
const discount := 10

func main:
    items := ["rice", "oil", "tea", "sugar"]
    prices := [180, 95, 120, 40]
    quantities := [2, 1, 1, 3]
    mut total := 0

    println "Receipt"
    println "-----------------------------"
    each (i, item) in items:
        amount := prices[i] * quantities[i]
        println item, "x", quantities[i], "=", amount
        total += amount
    end
    println "-----------------------------"
    println "Total:", total

    if total > threshold:
        discount amount := (total * discount) / 100
        println "Discount", discount, "percent:", discount amount
        println "Amount due:", total - discount amount
    else:
        println "Amount due:", total
        remaining := threshold - total
        println "Spend", remaining, "more to get 10 percent off"
    end
end
\end{salamcode}

\begin{salamout}
Receipt
-----------------------------
rice x 2 = 360
oil x 1 = 95
tea x 1 = 120
sugar x 3 = 120
-----------------------------
Total: 695
Discount 10 percent: 69
Amount due: 626
\end{salamout}

Change the quantities of rice and sugar both to 1 to see the \code{else}
branch as well. The variable \code{amount} is declared inside the loop and
without \code{mut}, because it receives a value once in each iteration and
keeps it; only \code{total} changes over the course of the loop.

\begin{tip}[An idea to take further]
Add a fourth array of booleans for ``discounted item'', and give those
items a separate 5 percent discount. Or find the most expensive line of the
receipt and print it at the end.
\end{tip}

\section{A one-month calendar}

\textbf{The problem.} We want to print the calendar of one month: seven
columns from Sunday to Saturday, with the days of the month in the right
places. For this we need to know how many days the month has and which day
of the week it starts on (0 for Sunday up to 6 for Saturday).

\textbf{The steps.} First we print the column headings. Then we build a
string and leave blank spaces for the days of the week that pass before the
first day of the month. Then we attach the days one by one, and every time we
reach Saturday we print the line and start a fresh one. To keep the columns
lined up, we print single-digit numbers with an extra space.

\program{A month calendar}
\begin{salamcode}
func two digits(n: int): str:
    if n < 10: ret " " + n end
    ret "" + n
end

func main:
    day count := 31
    // 0 Sunday ... 6 Saturday; this month starts on a Wednesday
    first day := 3

    println "Su  Mo  Tu  We  Th  Fr  Sa"
    mut line := ""
    mut column := 0

    // blank spaces before the first day
    repeat first day:
        line += "    "
        column++
    end

    repeat 1 to day count in day:
        line += two digits(day) + "  "
        column++
        if column == 7:
            println line
            line = ""
            column = 0
        end
    end
    if column > 0: println line end
end
\end{salamcode}

\begin{salamout}
Su  Mo  Tu  We  Th  Fr  Sa
             1   2   3   4
 5   6   7   8   9  10  11
12  13  14  15  16  17  18
19  20  21  22  23  24  25
26  27  28  29  30  31
\end{salamout}

The final condition \code{if column > 0} is there so that the unfinished
last week is printed too; without it, the days 26 to 31 would stay in
\code{line} and never be printed. This kind of ``unfinished business after
the loop'' is one of the places programmers most often forget.

\begin{tip}[An idea to take further]
With \code{match}, find the number of days of a month from its number, and
work out the first day of each month from the first day of the month before
it (the remainder of the sum on division by 7). Then you can print the
calendar of a whole year.
\end{tip}

\section{Analysing a sentence}

\textbf{The problem.} We have a sentence and want to know how many words it
has, which word is the longest, and how many times a particular word appears
in it.

\textbf{The steps.} We split the sentence at the spaces. With one loop over
the words we count them, find the longest with the ``maximum'' pattern, and
wherever a word equals the word we are looking for, we increase a counter.

\program{Analysing a sentence}
\begin{salamcode}
func main:
    sentence := "salam means peace and salam is for everyone"
    target := "salam"

    words := sentence.split(" ")
    mut longest := ""
    mut target count := 0

    repeat words.len() in i:
        word := words.get(i)
        if word.len() > longest.len():
            longest = word
        end
        if word == target:
            target count++
        end
    end

    println "Number of words:", words.len()
    println "Longest word:", longest
    println "\"" + target + "\" appears", target count, "times"
end
\end{salamcode}

\begin{salamout}
Number of words: 8
Longest word: everyone
"salam" appears 2 times
\end{salamout}

\code{longest} starts as the empty string, whose length is zero; so the
first word is sure to take its place, and after that only longer words do.
It is the ``maximum'' pattern again, this time for strings.

\begin{tip}[An idea to take further]
Find the shortest word too. Then build a new sentence in which every
occurrence of the target word is replaced with another word (say
``hello'').
\end{tip}

\section{What next?}

If you have come this far and understood the programs of this chapter, you
are a programmer; a beginner, perhaps, but a programmer. From here several
roads lie ahead:

\begin{itemize}
  \item \textbf{Write more.} Turn every small daily task that involves
        numbers or text into a program: the interest on a loan, a weekly
        study timetable, the water consumption of your home. Look again at
        the exercises at the ends of the chapters and write the ones you
        skipped.
  \item \textbf{Read the sample programs in the online editor.} You can
        understand all of them now. Some of them show things that were not
        in this book, such as \code{struct} for grouping related data, or
        \code{Vector} for lists that grow and shrink.
  \item \textbf{The next Salam books.} This book was the foundation. Salam
        has a great deal more to offer: a large library of ready-made
        tools, building web pages, working with files and networks, and much
        else. All of it is built on the same foundation you learned here.
\end{itemize}

Above all, do not be afraid of mistakes. Every error message you see is a
small lesson, and every program that finally works is a small victory. The
joy of programming lies in these small victories.

\begin{summary}
\begin{itemize}
  \item Break every problem into small steps, and write a function or a
        block for each clear step.
  \item The patterns you learned (the counter, the sum, the maximum,
        searching, building a string in a loop) come up again and again in
        real programs.
  \item After every loop, ask: is there unfinished business that still has
        to be done?
  \item A working program is not the finish line; make it better, shorter
        and more readable.
\end{itemize}
\end{summary}

\section*{Exercises}
\markright{Exercises}

\exercise Write a program that, for a loan with a given amount, annual
interest rate and term (in months), works out the simple monthly instalment
(the amount plus the total interest, divided by the number of months) and
prints a month-by-month table of the remaining debt.

\exercise Put the temperature of every day of a month into an array and
print a report that includes: the month's average, the hottest and coldest
days, the number of days above 30 degrees, and the longest run of
consecutive days above 30 degrees.

\exercise Put a list of employee names and their weekly working hours into
two parallel arrays. With a fixed hourly rate, work out each person's pay;
hours beyond 40 a week count as overtime at one and a half times the rate.
Print the total of all the pay as well.

\exercise Write a ``unit converter'': a function that takes a value and the
names of a source and a target unit (for example \code{"km"} and
\code{"m"}) and does the conversion with \code{match}. Support at least
five pairs of units.

\exercise Put a multi-line text (using \code{\textbackslash n}) into a
string, split it into lines with \code{split("\textbackslash n")}, and for
each line print its number and its number of words. At the end, say which
line has the most words.
